\uuid{GMkP}
\titre{Raisonnements}
\auteur{bodin ; croizat ; sacre}
\organisation{exo7}
\licence{CC BY-NC-SA 4.0}
\datecreate{2022-03-08}
\niveau{L1}
\module{Algèbre}
\chapitre{Logique, ensemble, raisonnement}
\sousChapitre{Absurde et contraposée}
\difficulte{2}
\theme{Raisonnements}
\exo7qcmid{hilisit.logique.050}

\contenu{
\begin{question}
\qclassification{100.03}
Téo et Théa jouent à un jeu de société. Téo est proche de la victoire ; il doit lancer un dé et Théa remarque avec raison que : "si Téo fait 4, alors il gagne le jeu". Quelles sont les affirmations certaines ?
\begin{answers}
	\bad{Si Téo fait 3, alors il n'aura pas gagné.}
	\bad{Si Téo gagne, c'est qu'il a fait 4.}
	\good{Si Téo ne gagne pas, c'est qu'il n'a pas fait 4.}
	\bad{Si Téo gagne fait 5, il perd.}
\end{answers}
\begin{explanations}
L'affirmation de Théa est "Dé $=4 \implies $ Téo gagne". Cette affirmation ne nous dit rien sur ce qui adviendra si le lancer de dé ne donne pas $4$ ! En revanche, la contraposée de l'affirmation nous informe que "Téo ne gagne pas $\implies$ Dé $\neq 4$".
\end{explanations}
\end{question}
}
